\documentclass[UTF8,11pt]{ctexart}
\usepackage{amsmath,amssymb,geometry}
\geometry{margin=22mm}
\begin{document}
\section*{P25: 多院区麻醉医师三阶段鲁棒人员规划SPP}
\paragraph{来源与范围}依据原论文PDF第1, 7, 8, 9, 10, 11, 12, 13, 14, 15, 16, 17页，(1)–(7)。
\subsection*{范围、集合与参数}
选定论文原始三阶段鲁棒人员规划SPP(7)。以手术日$t$为实例索引，分别优化各日的人员配置。$I$为麻醉医师集合，$S$为医院地点集合。

$u_{i,t}\in\{0,1\}$为当日可用，$v_{i,s}\in\{0,1\}$为常规分配可行地点；$r_{i,j,t}\in\{0,1\}$表示轮值顺序要求医师$i$先于$j$进入候召及被调用。$t_r,t_q\geq0$分别为常规和候召班可提供时长；$c_L,c_N,c_O,c_I\geq0$依次为候召调用、未被调用、单位加班、单位闲置成本。常规工资对当日方案为固定成本，不进入比较目标。二元排序约束取$M=1$已足够。

医院在术前三天报告的人员需求向量$D=(D_s)_{s\in S}$取自给定非空有限集合
\[
\mathcal D_t=\left\{D\in\mathbb Z_+^{|S|}:D_s=\bar D_{s,t}+f_s\widehat D_{s,t},\ |f_s|\leq1,\ \sum_{s\in S}|f_s|\leq\tau\right\}.
\]
$\bar D_{s,t}\geq0$是预测，$\widehat D_{s,t}\geq0$是最大偏差，$\tau\in[0,|S|]$控制联合不确定预算。给定$D$后，最终账单工时$B=(B_s)_{s\in S}$属于
\[
\mathcal B_t(D)=\prod_{s\in S}[\beta_{1,L}D_s+\beta_{0,s,t}^{L},\ \beta_{1,U}D_s+\beta_{0,s,t}^{U}],
\]
各区间非空、非负；$\beta$为外部历史数据拟合得到的参数。
\subsection*{决策、信息与可行域}
术前六周的$x_{i,s,t}\in\{0,1\}$表示常规地点分配，$p_{i,t}\in\{0,1\}$表示进入候召名单。其完整可行域$X_t$由下式定义：
\begin{align}
x_{i,s,t}&\leq u_{i,t}v_{i,s}&&i\in I,s\in S,\\
p_{i,t}&\leq\sum_{s\in S}u_{i,t}v_{i,s}&&i\in I,\\
p_{j,t}&\leq p_{i,t}+1-r_{i,j,t}&&i,j\in I,\\
p_{i,t}+\sum_{s\in S}x_{i,s,t}&\leq1&&i\in I,\\
x_{i,s,t},p_{i,t}&\in\{0,1\}&&i\in I,s\in S.
\end{align}
第二阶段在$D$揭示后、$B$揭示前选择$q_{i,s,t}(D)\in\{0,1\}$，表示从候召名单调用医师到地点$s$。它可依赖完整需求向量$D$，不得依赖尚未观察的$B$。给定$p$时可行域$P_t(p)$为
\begin{align}
q_{i,s,t}(D)&\leq p_{i,t}&&i\in I,s\in S,\\
q_{j,s,t}(D)&\leq q_{i,s,t}(D)+1-r_{i,j,t}&&i,j\in I,s\in S,\\
\sum_{s\in S}q_{i,s,t}(D)&\leq1&&i\in I,\\
q_{i,s,t}(D)&\in\{0,1\}&&i\in I,s\in S.
\end{align}
候召可行域由论文(5b)--(5e)的约束逐项定义，常规分配和候召安排分别使用各自的地点约束。
\subsection*{唯一完整鲁棒模型}
令$[z]^+=\max\{z,0\}$，提供的工时
$H_s(x,q)=\sum_{i\in I}(t_rx_{i,s,t}+t_qq_{i,s,t})$，以下两种费用只是表达式定义：
\begin{align*}
C_t(p,q)&=c_L\sum_{i\in I,s\in S}q_{i,s,t}+c_N\sum_{i\in I}\left(p_{i,t}-\sum_{s\in S}q_{i,s,t}\right),\\
U_t(B,x,q)&=\sum_{s\in S}\left(c_O[B_s-H_s(x,q)]^++c_I[H_s(x,q)-B_s]^+\right).
\end{align*}
\begin{equation}
\boxed{V_t=\min_{(x,p)\in X_t}\ \max_{D\in\mathcal D_t}\ \min_{q\in P_t(p)}\left\{C_t(p,q)+\max_{B\in\mathcal B_t(D)}U_t(B,x,q)\right\}.}
\end{equation}
上述嵌套顺序是一个不可拆分的三阶段应用模型：事前配置、需求揭示后调用、最终工时下结算。$X_t,P_t,\mathcal D_t,\mathcal B_t$分别定义各阶段决策域及完整需求不确定集。最优值对应上述三阶段嵌套问题。

\end{document}
